\chapter{Getting Access: Equity Markets}\label{ch:m1:getting-access-equities}

Three people have a strategy that works on paper, ten million dollars of
their own and their investors' money, and no way to send an order. No
exchange will take an order from them: exchanges deal with members. No
member will lend them its membership without controls, because the law makes
the member answerable for every order that passes through it. Nobody will
settle their trades until a clearing firm has agreed to stand behind them.
The data they need to decide what to trade is licensed, and the licence for a
machine that reads the data costs more than the one for a person. And the
price of each share they trade, to the hundredth of a cent, will depend on
how many they trade in a month. The first twenty-eight chapters described the
market as a mechanism. This chapter and the next describe it as a set of
contracts, which is how a firm meets it.

\section{Member or client?}

\begin{definition}[Executing and clearing brokers]\label{def:m1:getting-access-equities:brokers}
An \emph{executing broker}\index{executing broker} is the member through which
a firm's orders reach a venue, in the broker's name and under its pre-trade
risk controls (\cref{ch:m1:exchanges-brokers-venues}). A \emph{clearing
broker}\index{clearing broker} is the clearing-house member that settles the
firm's trades, carries its positions and finances them. A
\emph{give-up}\index{give-up (equities)} is a trade executed by one broker
and passed, by agreement, to another for clearing, so that a firm can execute
with many brokers and hold everything at one.
\end{definition}

There are three routes, in increasing order of cost, control and speed.

\textbf{Client of a broker.} The firm sends orders to a broker's algorithms or
smart router. No registration; the broker's fees, tiers and latency.

\textbf{Sponsored or \omterm{def:m1:exchanges-brokers-venues:dma}{direct market access}.} The firm's orders go to the
exchange through the broker's membership and, by law, through risk checks
that the broker controls; the firm chooses venue and order type itself, often
from its own servers in the exchange's data centre. It pays the broker per
share and receives the broker's exchange tier, in whole or in part.

\textbf{Registration and membership.} The firm becomes a \omterm{def:m1:the-sell-side:sell-side}{broker-dealer} and a
member of each exchange it needs. It faces the exchange directly, earns its
own tiers, may join liquidity-provider programmes, and takes on capital,
reporting, supervision and examination.

\begin{dated}{2026-09}{What registration costs in capital}\label{dat:m1:getting-access-equities:netcap}
Under the net capital rule a \omterm{def:m1:the-sell-side:sell-side}{broker-dealer} that carries customer accounts
must maintain net capital of at least \$250\,000; a dealer, which includes any
firm making more than ten trades a year for its own investment account, at
least \$100\,000; a broker that introduces its customers to a clearing firm,
\$50\,000; one that never holds customer funds or securities, \$5\,000. A
\omterm{def:m1:what-a-trading-firm-does:market-maker}{market maker} needs \$2\,500 for each security in which it makes a market
(\$1\,000 for those priced at \$5 or less), up to \$1 million. Aggregate
indebtedness may not exceed 1\,500\% of net capital. These are floors; the
binding constraint for a trading firm is the haircut on its positions.
\end{dated}

\begin{omfigure}
\begin{tikzpicture}[font=\small,
  box/.style={draw=omDef,rounded corners=2pt,minimum width=2.6cm,minimum height=0.95cm,align=center,fill=omDef!4},
  ext/.style={draw=omExo,rounded corners=2pt,minimum width=2.6cm,minimum height=0.95cm,align=center,fill=omExo!6},
  flow/.style={-{Stealth[length=2mm]},thick}]
\node[box] (f) at (0,0) {Trading firm};
\node[box] (eb) at (4.3,1.3) {Executing broker\\(risk checks)};
\node[ext] (ex) at (9.0,1.3) {Exchanges};
\node[box] (cb) at (4.3,-1.3) {Clearing broker\\(financing, margin)};
\node[ext] (ccp) at (9.0,-1.3) {Clearing house};
\draw[flow,omDef] (f.north east) -- node[above,sloped,font=\footnotesize]{orders} (eb.west);
\draw[flow,omDef] (eb) -- node[above,font=\footnotesize]{as member} (ex);
\draw[flow,omThm] (ex) -- node[right,font=\footnotesize]{trades} (ccp);
\draw[flow,omThm] (ccp) -- node[above,font=\footnotesize]{settles with} (cb);
\draw[flow,omThm] (cb.west) -- node[below=2pt,sloped,font=\footnotesize,pos=0.4]{positions, cash} (f.south east);
\draw[flow,omExo,dashed] (eb) -- node[right,font=\footnotesize]{give-up} (cb);
\end{tikzpicture}

\omcaption{A firm that is neither member nor clearing member. Orders go up
through the \omterm{def:m1:getting-access-equities:brokers}{executing broker}; trades come back down through the \omterm{def:m1:getting-access-equities:brokers}{clearing
broker}, which is the firm's real counterparty for money and
securities.}\label{fig:m1:getting-access-equities:chain}
\end{omfigure}

\section{The clearing arrangement}

The \omterm{def:m1:getting-access-equities:brokers}{clearing broker} is the contract that matters most and the one new firms
think about last. It decides how much leverage the firm has (its house
margin, not the regulatory minimum), what it pays to finance longs and
borrow shorts (\cref{ch:m1:financing}, \cref{ch:m1:stock-loan-and-short-selling}),
whether positions are margined as a portfolio, what it charges per share or
per ticket to clear, and the conditions under which it may raise margin or
liquidate without notice. A \omterm{def:m1:getting-access-equities:brokers}{clearing broker} takes the firm's credit risk
overnight; it will want to see the strategy's risk profile, the principals'
history, audited capital, and a minimum revenue commitment. A firm whose
\omterm{def:m1:getting-access-equities:brokers}{clearing broker} gives thirty days' notice has thirty days to find another or
close.

\section{Fee schedules and tiers}

\begin{definition}[Volume tier]\label{def:m1:getting-access-equities:tier}
A \emph{volume tier}\index{volume tier} is a line of an exchange's fee
schedule that grants a better rate to members whose monthly activity exceeds
a threshold, usually expressed as average daily added volume in proportion to
the total consolidated volume of the market. The rate of the tier reached
applies to all of the month's qualifying shares.
\end{definition}

\begin{dated}{2026-09}{One exchange's schedule}\label{dat:m1:getting-access-equities:bzx}
The equities fee schedule of one US maker-taker exchange effective 1
September 2026, for shares priced at \$1 or more: removing liquidity costs
\$0.0030 a share; adding displayed liquidity earns \$0.0016 at the standard
rate, and under the add-volume tiers \$0.0020 from 0.06\% of total
consolidated volume, \$0.0023 from 0.20\%, \$0.0028 from 0.25\% and, at the top
tier quoted here, \$0.0031 from 1.00\%. Non-displayed orders earn from nothing
to \$0.0008. The schedule runs to many pages of further tiers, by tape, by
order type and for combinations of activity across the group's exchanges.
\end{dated}

\begin{omfigure}
\begin{tikzpicture}
\begin{axis}[width=11cm,height=5cm,
  xlabel={average daily added volume (million shares)},ylabel={\$ thousand a month},
  tick label style={font=\small},label style={font=\small},
  xmin=0,xmax=40,ymin=0,ymax=2500,grid=major,grid style={gray!25},
  scaled y ticks=false,yticklabel style={/pgf/number format/fixed,/pgf/number format/1000 sep={\,}},
  table/col sep=comma]
\addplot[omDef,very thick,const plot] table[x=adav_m,y=monthly_k]{figdata/markets-1/29-getting-access-equities/tiers.csv};
\node[font=\footnotesize,anchor=west] at (axis cs:7.6,230) {0.06\%};
\node[font=\footnotesize,anchor=west] at (axis cs:24.4,1030) {0.20\%};
\node[font=\footnotesize,anchor=west] at (axis cs:30.4,1560) {0.25\%};
\end{axis}
\end{tikzpicture}

\omcaption{Monthly rebates under the schedule of
\cref{dat:m1:getting-access-equities:bzx}, with total consolidated volume
assumed at 12 billion shares a day and 21 trading days. Because a tier's rate
applies to the whole month, each threshold is a cliff: the share that crosses
0.20\% is worth \$155\,000. Data: the chapter's build.}\label{fig:m1:getting-access-equities:tiers}
\end{omfigure}

Three consequences follow. \emph{Scale is an edge}: two firms with the same
strategy earn different amounts per share, and the larger can quote where the
smaller cannot. \emph{Aggregation is a business}: a broker whose clients
together reach a high tier can pass most of that rate to each of them and
keep the rest, which is often a better deal for a small firm than its own
membership. \emph{Volume near a threshold has a shadow price}: late in the
month a firm just under a tier will trade shares at a loss to get there, and
its behaviour changes on the day it arrives.

\begin{proposition}[When padding volume pays]\label{prop:m1:getting-access-equities:padding}
A member adds $v$ shares a day naturally, earning a rebate $\rho_0$ a share.
The next tier, at $V^* > v$, pays $\rho_1 > \rho_0$. Extra volume can be had at
a loss of $\ell$ a share before rebates, $\ell > \rho_1$. Padding up to $V^*$
is profitable if and only if
\[
v \;>\; V^*\,\frac{\ell - \rho_1}{\ell - \rho_0}.
\]
\end{proposition}
\begin{proof}
Padding changes the daily result by $V^*\rho_1 - (V^*-v)\ell - v\rho_0$, which
is positive exactly when $v(\ell-\rho_0) > V^*(\ell-\rho_1)$.
\end{proof}

\begin{omfigure}
\begin{tikzpicture}
\begin{axis}[width=10.5cm,height=4.8cm,ybar,bar width=12pt,
  ylabel={cents per 100 shares},
  symbolic x coords={maker-taker,inverted,flat-fee},xtick=data,
  tick label style={font=\small},label style={font=\small},xticklabel style={font=\small\strut},
  ymin=-25,ymax=42,ymajorgrids,grid style={gray!25},enlarge x limits=0.22,
  nodes near coords,nodes near coords style={font=\footnotesize,fill=white,inner sep=1pt,/pgf/number format/.cd,fixed,precision=0},
  legend columns=2,legend style={font=\footnotesize,at={(0.5,-0.2)},anchor=north,draw=none,fill=none,/tikz/every even column/.append style={column sep=8pt}},
  table/col sep=comma]
\addplot[fill=omDef!60,draw=omDef] table[x=venue,y=passive_mils]{figdata/markets-1/29-getting-access-equities/venues.csv};
\addplot[fill=omThm!60,draw=omThm] table[x=venue,y=aggressive_mils]{figdata/markets-1/29-getting-access-equities/venues.csv};
\legend{resting order executed,marketable order}
\end{axis}
\end{tikzpicture}

\omcaption{All-in cost per 100 shares on three kinds of venue for a member
adding 9 million shares a day: exchange fee or rebate, plus 2 cents of
clearing and 1 cent of regulatory fees. Negative is income. The maker-taker
rates are the published ones; the inverted and flat venues and the other two
costs are illustrative. Data: the tutorial.}\label{fig:m1:getting-access-equities:venues}
\end{omfigure}

The sign of the fee decides where an order should rest. On the \omterm{def:m1:us-equity-market-structure:makertaker}{inverted
venue} a resting order \emph{pays}, so the queue is short and fills come
first; on the maker-taker venue it is paid, so the queue is long and the
order is filled last, when the price is about to move through it. The
difference in fees, three tenths of a cent, is the market's price for queue
priority, and a router that ignores it
(\cref{ch:m1:us-equity-market-structure}) misprices every passive order.

\section{Liquidity-provider programmes}

\begin{definition}[Designated market maker and supplemental liquidity provider]\label{def:m1:getting-access-equities:dmm}
A \emph{designated market maker}\index{designated market maker} is the
member assigned to a listed security with obligations to maintain a fair and
orderly market in it, including its opening and closing auctions, in return
for information and allocation privileges. A \emph{supplemental liquidity
provider}\index{supplemental liquidity provider} is a member that commits to
quoting assigned securities competitively for a minimum share of the day, in
return for enhanced rebates, without the designated maker's other
duties.
\end{definition}

\begin{dated}{2026-09}{A quoting requirement}\label{dat:m1:getting-access-equities:slp}
On the New York Stock Exchange a \omterm{def:m1:getting-access-equities:dmm}{supplemental liquidity provider} must
maintain a bid or an offer at the national best bid or offer in each assigned
security for at least 10\% of the trading day; it may be a proprietary trading
unit of a member or a registered \omterm{def:m1:what-a-trading-firm-does:market-maker}{market maker}, and the two are not
aggregated for the requirement.
\end{dated}

Programmes of this kind exist on every exchange under different names.
They are contracts: obligations measured by the exchange (time at the inside,
quoted size, spread), against rates not available otherwise. A firm already
quoting at the inside most of the day should be in the programme; a firm that
would have to change its behaviour to qualify is selling an option on its
quotes for the rebate, and should price it.

\section{Data licences}

\begin{definition}[Non-display fee]\label{def:m1:getting-access-equities:nondisplay}
A \emph{non-display fee}\index{non-display fee} is the charge an exchange
makes for the use of its market data by a machine (an algorithm, a router, a
risk system) as opposed to a person looking at a screen. It is levied per
firm or per platform, by category of use, on top of the fee for receiving the
data.
\end{definition}

A trading firm's data bill has four layers: access (the feed itself and the
ports and cross-connects it arrives on), redistribution (if data goes to
other entities, even affiliates), display (per user) and non-display (per
category of automated use). Exchanges audit. The bill is per exchange, and a
strategy that needs the \omterm{def:m1:us-equity-market-structure:sip}{direct feeds} of a dozen exchanges
(\cref{ch:m1:reading-market-data}) pays a dozen of them before its first
trade. Market data is a significant line of an exchange group's revenue
(\cref{ch:m1:exchanges-brokers-venues}) and, for a small firm, of its costs.

\section{The negotiation}

Nothing on a published exchange schedule is negotiable: exchange fees are
filed with the regulator and apply to every member that meets their terms,
which is why the schedules are long. Everything else is: the broker's pass-through of its tier; the
clearing fee per share and its minimum; financing spreads and borrow rates;
the portion of the \omterm{def:m1:getting-access-equities:brokers}{clearing broker}'s stock-loan revenue on the firm's longs
that comes back; cross-connect and hosting prices at third-party data
centres; vendors' data and software. What the firm brings to each
negotiation is volume, balances and the credible ability to go elsewhere. It
should therefore know its own numbers better than the counterparty does:
shares per day by venue and by add or remove; average balances, long and
short; hard-to-borrow usage; messages per second. The builds of this book
produce most of them.

\section{Tutorial: net cost per share on three venues}\label{tut:m1:getting-access-equities:cost}

\begin{tutorial}
\textbf{Goal.} Encode a tiered schedule, find a member's tier, compute the
month's bill, measure the cliff at a threshold, and compare the all-in cost
of a passive and an aggressive share on three kinds of venue. \textbf{End
state:} the two data figures of this chapter.
\begin{tutsteps}
\item \textbf{A schedule} is a base rate, a remove rate and a list of tiers.
\omcode{code/firm/feesched/firm_feesched.py}{18}{41}{Tier lookup and the daily bill.}\label{lst:m1:getting-access-equities:schedule}
\item \textbf{Padding.} The better rate applies to everything; the padded
shares lose money before the rebate.
\omcode{code/firm/feesched/firm_feesched.py}{44}{54}{Is it worth adding volume to reach the next tier?}\label{lst:m1:getting-access-equities:padding}
\item \textbf{All-in.} Exchange, clearing, regulatory.
\end{tutsteps}
\textbf{What to change next.} Add a second condition to a tier (a minimum
share of volume \emph{removed} as well) and a cross-exchange tier that counts
activity on a sister exchange; then find the cheapest allocation of a fixed
order flow across the group's exchanges.
\end{tutorial}

\section{Build: the fee-schedule engine}\label{bld:m1:getting-access-equities:feesched}

\begin{build}
\bfield{Purpose} The miniature firm's backtester (One Quant Book 7) and
router (Book 11) need the marginal fee of every order, which depends on the
month's volume so far; its finance function needs the bill.
\bfield{Interface} \texttt{Tier(name, min\_adav\_share, add\_rate)};
\texttt{Schedule(venue, base\_add\_rate, remove\_rate, tiers)} with
\texttt{tier\_for}, \texttt{add\_rate}, \texttt{daily\_bill},
\texttt{next\_tier}; \texttt{extra\_volume\_worth\_adding(schedule, adav, tcv,
loss)}; \texttt{breakeven\_natural\_volume(current, next, threshold, loss)};
\texttt{all\_in\_per\_share(\ldots)}.
\bfield{Rules} Rates per share, signed: positive is a cost. The best
qualifying tier applies to all added shares of the period. Thresholds are
shares of total consolidated volume, so the engine takes that volume as an
input and a forecast of it for the current month. Schedules are data with an
effective date, never code: they change monthly.
\bfield{Acceptance tests} \texttt{code/firm/feesched/tests/}: tier lookup at
and around thresholds; the cliff at 0.20\%; padding that loses and padding
that pays, with the closed-form break-even; all-in costs.
\bfield{Stretch} Parse a real schedule's tier table from its published
page into the data format, with a diff against last month's.
\end{build}

\begin{omsources}
\item Cboe BZX Equities, \emph{Fee Schedule}, effective 1 September 2026.
\item 17 CFR 240.15c3-1 (net capital requirements for brokers or dealers).
\item New York Stock Exchange, Rule 107B (supplemental liquidity providers)
as described in the exchange's fee filings with the SEC.
\end{omsources}

\section{Exercises}

\begin{exercise}[$\star$]\label{exo:m1:getting-access-equities:1}
A member adds 9 million shares a day when total consolidated volume is 12
billion. Find its tier under \cref{dat:m1:getting-access-equities:bzx}, its
daily rebate and its rebate for a 21-day month.
\end{exercise}

\begin{exercise}[$\star$]\label{exo:m1:getting-access-equities:2}
Compute the daily rebate at 23.9 and at 24.0 million shares added. What is the
hundred-thousandth share worth?
\end{exercise}

\begin{exercise}[$\star$]\label{exo:m1:getting-access-equities:3}
With clearing at 2 cents and regulatory fees at 1 cent per 100 shares, give
the all-in cost per 100 shares of a passive and of an aggressive execution
for the member of exercise 1, and of a round trip that enters passively and
exits aggressively.
\end{exercise}

\begin{exercise}[$\star\star$]\label{exo:m1:getting-access-equities:4}
A firm trades for its own account only and makes markets in 300 securities
priced above \$5. Using \cref{dat:m1:getting-access-equities:netcap}, give its
minimum net capital. With \$2 million of net capital, what is its maximum
aggregate indebtedness? Why is neither number the one that limits its
trading?
\end{exercise}

\begin{exercise}[$\star\star$]\label{exo:m1:getting-access-equities:5}
A \omterm{def:m1:getting-access-equities:dmm}{supplemental liquidity provider} must be at the inside for 10\% of the day in
each assigned security. How many minutes is that? A firm is at the inside
bid 6\% of the day and at the inside offer 7\%, never both at once. Does it
qualify? What if the two always coincide?
\end{exercise}

\begin{exercise}[$\star\star$]\label{exo:m1:getting-access-equities:6}
On the \omterm{def:m1:us-equity-market-structure:makertaker}{inverted venue} of \cref{fig:m1:getting-access-equities:venues} a
resting order costs 13 cents per 100 shares all-in, against an income of 17 on
the maker-taker venue. Give two reasons why a \omterm{def:m1:what-a-trading-firm-does:market-maker}{market maker} might still rest
orders there.
\end{exercise}

\begin{exercise}[$\star\star\star$]\label{exo:m1:getting-access-equities:7}
\emph{Coding.} With \texttt{breakeven\_natural\_volume} find the natural volume
above which padding to the 0.20\% tier pays, for a loss of 27 cents and of 25
cents per 100 padded shares. Check each with
\texttt{extra\_volume\_worth\_adding} just above and just below.
\end{exercise}

\begin{exercise}[$\star\star\star$]\label{exo:m1:getting-access-equities:8}
\emph{Find the flaw.} A business plan: ``Our passive strategy earns 10 cents
per 100 shares before fees. The exchange pays 31 cents per 100 shares for
adding liquidity. Net: 41 cents per 100 shares on 2 million shares a day.''
Find two errors and recompute.
\end{exercise}

\section{Problem: Reaching the Tier}

\begin{problem}[{Weekend problem --- the last week of the month}]\label{pb:m1:getting-access-equities:1}
A member adds 20 million shares a day naturally under the schedule of
\cref{dat:m1:getting-access-equities:bzx}; total consolidated volume is 12
billion a day; the month has 21 days. Extra passive volume can be generated
by quoting more aggressively, at a loss of 27 cents per 100 shares before
rebates.

\medskip\noindent\textbf{Part I --- The cliff.}
\begin{enumerate}
\item Give the member's tier and monthly rebate.
\item Give the volume needed for the next tier and the extra volume a day.
\item Give the daily gain or loss from padding up to it.
\item Give the natural volume at which padding breaks even.
\item The same for the 0.25\% tier, starting from 24 million. Does padding from
24 to 30 million pay?
\end{enumerate}

\medskip\noindent\textbf{Part II --- Member or client?}
A broker whose clients together reach the 0.25\% tier offers to pass through
80\% of its rebate. Membership costs the firm \$25\,000 a month in fixed fees
(membership, ports, connectivity, compliance).
\begin{enumerate}[resume]
\item Give the client's effective rebate per share.
\item At 20 million shares a day, compare the two routes per month.
\item At 30 million, compare again.
\item What else does the firm give up as a client?
\item What else does it take on as a member?
\end{enumerate}

\medskip\noindent\textbf{Part III --- The month as it happens.}
\begin{enumerate}[resume]
\item On day 15 the firm has averaged 23 million shares a day. What average
does it need over the last six days to reach 24 million for the month?
\item Total consolidated volume turns out to be 13 billion, not 12. What is
the threshold now, and what does that do to the plan?
\item How should the backtester charge fees to a strategy whose volume
decides the firm's tier?
\item Why do exchanges design schedules with cliffs?
\end{enumerate}

\medskip\noindent\textbf{Part IV --- Judgement.}
\begin{enumerate}[resume]
\item Is volume traded only to reach a tier harmful to anyone?
\item A competitor ten times larger earns the top tier. Quantify its
advantage per share and say where it shows up in the market.
\item What would a regulator's cap on rebates do to this chapter?
\item List what the firm should measure before negotiating with a broker.
\item State the \emph{named result}: the natural volume at which the better
tier pays for itself.
\item In one sentence: what does a fee schedule reward?
\end{enumerate}
\end{problem}

\section{Interview questions}

\begin{interviewq}[$\star$ \iqroles{trader, developer, researcher}]\label{iq:m1:getting-access-equities:1}
What is the difference between an \omterm{def:m1:getting-access-equities:brokers}{executing broker} and a \omterm{def:m1:getting-access-equities:brokers}{clearing broker}?
\end{interviewq}

\begin{interviewq}[$\star$ \iqroles{trader, researcher}]\label{iq:m1:getting-access-equities:2}
What is a maker-taker fee schedule, and what is a \omterm{def:m1:getting-access-equities:tier}{volume tier}?
\end{interviewq}

\begin{interviewq}[$\star\star$ \iqroles{trader, researcher}]\label{iq:m1:getting-access-equities:3}
Why would anyone rest an order on a venue that charges for adding
liquidity?
\end{interviewq}

\begin{interviewq}[$\star\star$ \iqroles{researcher, mle}]\label{iq:m1:getting-access-equities:4}
How should exchange fees enter a backtest?
\end{interviewq}

\begin{interviewq}[$\star\star$ \iqroles{trader, bank}]\label{iq:m1:getting-access-equities:5}
You are starting a \omterm{def:m1:what-a-trading-firm-does:prop}{proprietary trading firm} in US equities. Walk me through
the decisions about access.
\end{interviewq}

\begin{interviewq}[$\star\star\star$ \iqroles{trader, researcher}]\label{iq:m1:getting-access-equities:6}
It is the 25th of the month and you are 4\% short of a tier. What do you do,
and what do you not do?
\end{interviewq}
